% ═══════════════════════════════════════════════════════════════════════════════
%  CAPÍTULO — Iterator
% ═══════════════════════════════════════════════════════════════════════════════

\chapter{Iterator}

% ─── Situación Problema ──────────────────────────────────────────────────────
\section{Situación Problema}

Airbnb mantiene las propiedades disponibles en distintas estructuras
de datos internas según el contexto: una lista ordenada por relevancia
para los resultados de búsqueda, un árbol geoespacial para la vista
de mapa y un grafo de recomendaciones para la sección «también te
podría gustar». Los distintos módulos de la plataforma —exportación de
datos, motor de recomendaciones, generador de reportes— necesitan
recorrer estas colecciones de propiedades de forma uniforme, sin acoplarse
a su estructura interna.

Si cada módulo accede directamente a la representación interna de la
colección, cualquier cambio en la estructura de datos obligaría a
actualizar todos los módulos que la recorren, creando un acoplamiento
innecesario y frágil.

% ─── Solución ────────────────────────────────────────────────────────────────
\section{Solución}

El patrón Iterator define la interfaz \texttt{IteradorPropiedades} con
los métodos \texttt{siguiente()} y \texttt{hayMas()}. Cada estructura
de datos proporciona su propio iterador concreto que sabe cómo recorrerla
sin exponer su representación interna. Los módulos clientes utilizan
únicamente la interfaz del iterador y son agnósticos a si están
traversando una lista, un árbol o un grafo.

Cuando el equipo de backend decide cambiar la estructura interna de
almacenamiento de propiedades, basta con actualizar el iterador
correspondiente. El código de los módulos consumidores no sufre ningún
cambio, manteniéndose el principio de encapsulamiento de las colecciones.

% ─── Diagrama de clases ─────────────────────────────────────────────────────
\begin{figure}[H]
    \centering
    % \includegraphics[width=\textwidth]{imagenes/abstract_factory_diagrama.png}
    \caption{Diagrama de clases del patrón Abstract Factory}
\end{figure}


% ─── Roles de las Clases ────────────────────────────────────────────────────
\section{Roles de las Clases}

% Explicar la responsabilidad de cada clase/interfaz

\begin{itemize}
    \item \textbf{AbstractFactory:}
    \item \textbf{ConcreteFactory:}
    \item \textbf{AbstractProduct:}
    \item \textbf{ConcreteProduct:}
    \item \textbf{Client:}
\end{itemize}


% ─── Main ───────────────────────────────────────────────────────────────────
\section{Main}

% Explicar qué realiza el programa principal

\begin{lstlisting}[language=Java, caption=NombreClase]
 
\end{lstlisting}


% ─── Salida del Main ────────────────────────────────────────────────────────
\section{Salida del Main}

\begin{lstlisting}[language=Java, caption=NombreClase]
 
\end{lstlisting}


% ─── Clases ─────────────────────────────────────────────────────────────────
\section{Clases}

% ─── Clase 1 ────────────────────────────────────────────────────────────────
\subsection{NombreClase}

\begin{lstlisting}[language=Java, caption=NombreClase]
 
\end{lstlisting}


% ─── Clase 2 ────────────────────────────────────────────────────────────────
\subsection{NombreClase}

\begin{lstlisting}[language=Java, caption=NombreClase]
 
\end{lstlisting}


% ─── Clase 3 ────────────────────────────────────────────────────────────────
\subsection{NombreClase}

\begin{lstlisting}[language=Java, caption=NombreClase]
 
\end{lstlisting}


% ─── Conclusiones ───────────────────────────────────────────────────────────
\section{Conclusiones}

% Explicar:
% - Ventajas del patrón
% - Cuándo usarlo
% - Beneficios obtenidos
% - Posibles desventajas
% - Relación con principios SOLID